% ============================================================================
%  Chapter 15 -- Memory                                          (Lecture 15)
% ============================================================================
\chapter{Memory}
\label{ch:memory}

\epigraphline{\ldots computers and the programs will start thinking and the
people will stop.}{Dr.\ Walter Gibbs, \emph{Tron}, Steven Lisberger, 1982}

\begin{chapgoals}
By the end of this chapter you will be able to
\begin{itemize}
\item classify memories as volatile or persistent, on-line or off-line, and place them in the memory hierarchy;
\item describe how SRAM and DRAM store a bit, and why they are used for different jobs;
\item explain what words, addresses and byte-addressable memory are;
\item tell little-endian from big-endian order, and say which one \riscv uses;
\item describe step by step how the processor reads and writes a word;
\item explain memory-mapped input/output and what DMA is for.
\end{itemize}
\end{chapgoals}

% ----------------------------------------------------------------------------
\section{Warm-up on the machine}

\asmfile[caption={The same memory read as a word and as bytes.},label=lst:mem:warm]{code/l15/c03.S}

In Listing~\ref{lst:mem:warm} a word of thirty-two bits occupies four
consecutive addresses: memory is addressed byte by byte. The byte at the
lowest address is \hex{78}, the least significant: the order is little-endian.
\ins{lw} fetches four bytes in one operation, \ins{lbu} fetches one.

\begin{prgtry}
Run it in \prgqt and open the memory view at the address in \reg{t0}: the
bytes appear as \texttt{78 56 34 12}. The byte order of
Chapter~\ref{ch:information} has become an experiment.
\end{prgtry}

What is a memory made of, and what happens between the processor and the
memory when \ins{lw} runs?

% ----------------------------------------------------------------------------
\section{What is a memory?}

\begin{definition}{Memory}
The \term{memory} is the element of a computer, or of one of its subsystems,
dedicated to the persistence of data and of the instructions of programs.
\end{definition}

Memories are the components able to store information: data, programs and
results. They are classified along two axes. A memory is \term{volatile} if
the information is kept only as long as the computer is powered, and
\term{persistent} if it is kept even when the computer is switched off. It is
\term{on-line} if the data are always accessible, and \term{off-line} if the
medium must be mounted before the data can be reached.

\subsection{The memory hierarchy}

No technology is at the same time fast, large and cheap. A computer
therefore uses several, arranged in a \term{hierarchy}
(Figure~\ref{fig:mem:hier}): going up, memories are faster and more expensive
per bit; going down, larger and cheaper per bit. Typical access times are
below a nanosecond to a few nanoseconds for registers and cache, tens of
nanoseconds for DRAM, tens of microseconds for solid-state disks,
milliseconds for magnetic disks, seconds to minutes for tape.

\begin{figure}[htb]
\centering
\begin{tikzpicture}
  \foreach \i/\n/\d/\c/\w in {0/registers/{inside the processor}/prgred/30, 1/cache/{SRAM, on the processor chip}/prgamber/46, 2/{main memory}/DRAM/prgaccent/62, 3/{solid-state and magnetic disks}/{persistent, on-line}/prgsteel/78, 4/{tapes and removable media}/{persistent, off-line}/prgviolet/94}{
    \node[blk,fill=\c,draw=\c,text=white,minimum width=\w mm,minimum height=9mm] at (0,-\i*1.05) {\textbf{\n}\\[-1pt]\scriptsize \d};
  }
  \draw[->,very thick,prgink] (5.6,-4.4) -- (5.6,0.3) node[above,lbl,align=center]{faster, more\\expensive per bit};
  \draw[->,very thick,prgink] (-5.6,0.3) -- (-5.6,-4.4) node[below,lbl,align=center]{larger,\\cheaper per bit};
\end{tikzpicture}
\caption{The memory hierarchy.}
\label{fig:mem:hier}
\end{figure}

% ----------------------------------------------------------------------------
\section{A little history: rings of ferrite}

\subsection{Magnetic-core memory}

In the first computers, before semiconductor memories became widespread, a
bit was the direction of the magnetic field of a small ring of ferrite, a
magnetic ceramic (Figure~\ref{fig:mem:core}). A core memory is static: it
keeps its contents without power. The first research was done by An Wang and
Way-Dong Woo at Harvard in 1949; at the beginning of the 1950s Jay Forrester
at MIT built a core memory for Whirlwind, a computer meant as a flight
simulator.

\begin{figure}[htb]
\centering
\begin{subfigure}[b]{0.3\textwidth}\centering\includegraphics[width=\textwidth,height=38mm,keepaspectratio]{core_memory.jpg}\caption{Ferrite cores and their wires}\label{fig:mem:core}\end{subfigure}\hfill
\begin{subfigure}[b]{0.34\textwidth}\centering\includegraphics[width=\textwidth,height=38mm,keepaspectratio]{ram_modules.jpg}\caption{Memory modules of several generations}\label{fig:mem:modules}\end{subfigure}\hfill
\begin{subfigure}[b]{0.3\textwidth}\centering\includegraphics[width=\textwidth,height=38mm,keepaspectratio]{ram_on_board.jpg}\caption{Modules in their slots}\label{fig:mem:slot}\end{subfigure}
\caption{Memory, yesterday and today.}
\end{figure}

\subsection{Reading and writing a core}

Each core is threaded by wires. Two address lines, X and Y, select one core:
only the core at their crossing receives enough current to change.

\emph{Reading is destructive.} The selected bit is set to 0; if it contained
1, a pulse appears on a third wire, the \emph{sense} line, and otherwise
nothing happens. The datum has been read and lost. Wang's solution was that
every read is followed by a write that stores the datum again.

\emph{Writing.} All the bits start at 0. X and Y select the core, which is
set to 1; if the bit must be 0, an \emph{inhibit} line prevents the change.
(Sense and inhibit share one wire.) Forrester's contribution was the
arrangement that manages a large number of cores with few connections.

The destructive read followed by a rewrite has returned today in DRAM.

% ----------------------------------------------------------------------------
\section{Memory today}

\subsection{SRAM and DRAM}

\begin{definition}{SRAM and DRAM}
In a \term{static RAM} (SRAM) each bit is held by a bistable circuit, as in a
flip-flop. In a \term{dynamic RAM} (DRAM) each bit is a charge in a tiny
capacitor, behind one transistor.
\end{definition}

The cells of an SRAM are arranged in a matrix and reached by row and column;
each takes six transistors (Chapter~\ref{ch:electronics}) and keeps its
content as long as it is powered. SRAM is fast and expensive, and is used for
cache memories; it is also the main memory of a microcontroller such as the
ESP32-C6, which has 512~KB of SRAM and no DRAM at all.

The capacitor of a DRAM cell discharges, and its content must be
\term{refreshed} periodically by a circuit inside the memory. DRAM is cheap
and dense, and is the main memory of laptops, telephones and servers. An
\emph{SDRAM} is a synchronous DRAM, whose access is governed by a clock; DDR4
and DDR5 are generations of it.

\subsection{Other memories}

A \emph{ferroelectric RAM} (FeRAM) keeps its state without power.
\emph{Flash} memory is persistent and is erased in blocks: it is the memory of
USB keys, of telephones and of solid-state disks, treated in the next chapter.
The program of a microcontroller lives in flash and its variables in SRAM.

Memory modules have changed shape and size through the years
(Figure~\ref{fig:mem:modules}): SIMM, SDRAM PC100 and PC133, Rambus and DDR
held tens to hundreds of megabytes each; today's DDR5 modules hold 16 to
64~GB.

\subsection{How to recognise the memory}

On the motherboard of a desktop or of a server the memory is in rows of long
slots next to the processor, which hold the modules
(Figure~\ref{fig:mem:slot}). On a laptop or a telephone it is in chips
soldered close to the processor. On the \prg kit there is nothing to see: the
SRAM is inside the ESP32-C6 itself, and the flash is a small chip beside it.

Two parameters characterise a memory: its \term{size}, the number of cells,
which today is gigabytes on a PC and kilobytes on a microcontroller; and its
\term{access time}, the delay after which the datum is available, tens of
nanoseconds for DRAM. For comparison, at 160~MHz one clock period of the
ESP32-C6 is 6.25~ns.

% ----------------------------------------------------------------------------
\section{The main memory}

\subsection{Words and addresses}

\begin{definition}{Main memory}
The \term{main memory} is an array of \term{words}; a word is a string of $m$
bits, which the processor accesses with a single read or write operation.
Each word is identified by an \term{address}: an integer between 0 and $N-1$,
represented on $n$ bits, with $N=2^{n}$.
\end{definition}

The address is a code word, and the memory decodes it: it is the decoder of
Chapter~\ref{ch:combinational} and the small memory of
Chapter~\ref{ch:registers}, scaled up. Thirty-two address bits give $2^{32}$
addresses, 4~GiB. A memory is called \term{random access} (RAM) when the
access time does not depend on the position of the datum. Processors access
memory with a parallelism equal to the length of their word: 16, 32 or 64
bits.

\subsection{Byte-addressable memory}

In a \term{byte-addressable} memory the byte is the smallest unit that has an
address (Figure~\ref{fig:mem:bytes}). A word of 32 bits occupies four
consecutive addresses, and its address is that of its first byte.
Consecutive words are 4 addresses apart: this is why the program counter
advances by 4, and why the index of an array of words is multiplied by 4
(Chapter~\ref{ch:information}).

\begin{figure}[htb]
\centering\small
\begin{tabular}{@{}r@{\quad}*{8}{|>{\centering\arraybackslash\ttfamily}p{9mm}}|@{}}
\multicolumn{1}{r}{\footnotesize address} & \multicolumn{1}{c}{\scriptsize $A$} & \multicolumn{1}{c}{\scriptsize $A{+}1$} & \multicolumn{1}{c}{\scriptsize $A{+}2$} & \multicolumn{1}{c}{\scriptsize $A{+}3$} & \multicolumn{1}{c}{\scriptsize $A{+}4$} & \multicolumn{1}{c}{\scriptsize $A{+}5$} & \multicolumn{1}{c}{\scriptsize $A{+}6$} & \multicolumn{1}{c}{}\\\cline{2-9}
 & \cellcolor{prgaccent!20}78 & \cellcolor{prgaccent!20}56 & \cellcolor{prgaccent!20}34 & \cellcolor{prgaccent!20}12 & \cellcolor{prgamber!25}48 & \cellcolor{prgamber!25}45 & \cellcolor{prgamber!25}4C & \ldots\\\cline{2-9}
\multicolumn{1}{r}{} & \multicolumn{4}{c}{\footnotesize one word: \texttt{.word 0x12345678}} & \multicolumn{3}{c}{\footnotesize three bytes: \texttt{"HEL\ldots"}} & \multicolumn{1}{c}{}\\
\end{tabular}
\caption{A byte-addressable memory: each cell is one byte.}
\label{fig:mem:bytes}
\end{figure}

\subsection{Big-endian and little-endian}

The two names say how the bytes of a word are arranged in memory
(Figure~\ref{fig:mem:endian}). In \term{little-endian} order the least
significant byte goes at the lowest address: it is the order of \riscv, of
x86 and of ARM. In \term{big-endian} order the most significant byte goes at
the lowest address: it is the order of the Motorola~68000, the reference
processor of the previous editions of this course, and of the protocols of
the Internet.

\begin{figure}[htb]
\centering\small
\begin{tabular}{@{}r@{\quad}*{4}{|>{\centering\arraybackslash\ttfamily}p{13mm}}|@{}}
\multicolumn{1}{r}{} & \multicolumn{1}{c}{\scriptsize 1000} & \multicolumn{1}{c}{\scriptsize 1001} & \multicolumn{1}{c}{\scriptsize 1002} & \multicolumn{1}{c}{\scriptsize 1003}\\\cline{2-5}
\textbf{little-endian}: \riscv, x86, ARM & \cellcolor{prggreen!25}FF & \cellcolor{prgviolet!20}E6 & \cellcolor{prgamber!25}A3 & \cellcolor{prgaccent!20}10\\\cline{2-5}
\multicolumn{5}{c}{}\\[-6pt]\cline{2-5}
\textbf{big-endian}: Motorola 68000, network order & \cellcolor{prgaccent!20}10 & \cellcolor{prgamber!25}A3 & \cellcolor{prgviolet!20}E6 & \cellcolor{prggreen!25}FF\\\cline{2-5}
\end{tabular}
\caption{The word \hex{10A3E6FF} stored at address 1000.}
\label{fig:mem:endian}
\end{figure}

\subsection{Aligned and unaligned words}

A word that starts at an address which is a multiple of its size is
\term{aligned}: multiples of 4 for words of 32 bits. An aligned word is
fetched with a single memory access; an unaligned word straddles two memory
words, and two accesses are needed. The \riscv specification leaves the
treatment of misaligned accesses to the implementation: a load or a store to a
misaligned address is, depending on the processor, done more slowly or refused
with an exception. Correct programs keep their data aligned. The assembler
aligns the data declared with \code{.word}, and the directive
\code{.balign 4} forces alignment after a string or a byte
(Listing~\ref{lst:mem:align}). Instructions must always be aligned.

\asmfile[caption={Forcing the alignment of a word.},label=lst:mem:align,numbers=none]{code/l15/c04.S}

\subsection{Parameters of a memory}

The \term{capacity} of a memory is the number $N$ of its registers and the
length in bits of each: 1~MB is $2^{20}$ registers of 8 bits; ``16 kW of 32
bits'' is $16\times 2^{10}$ registers of 32 bits. The access time is the time
needed to carry out a read or a write: the time to select the location, plus
the time for the read or the write itself.

\begin{prgmachine}
The ESP32-C6 has 512~KB of SRAM for data, 320~KB of ROM, and an external
flash of some megabytes for programs. The menu of \prg shows what is left
for the programmer: about 98~KB of RAM and 509~KB of flash for cartridges.
The main memory of the chip is that of eight Commodore~64s, or one
thirty-thousandth of that of a laptop.
\end{prgmachine}

% ----------------------------------------------------------------------------
\section{Processor and memory}

\subsection{The interface}

The processor and the memory talk through three buses and two registers
(Figure~\ref{fig:mem:iface}). The \term{memory address register} (MA) holds
the address and drives the \term{address bus}. The \term{memory buffer
register} (MB) holds the datum that travels, in either direction, on the
\term{data bus}. The control unit drives the \term{control lines}, among
which are the read enable and the write enable. MA and MB are internal
registers of the model of processor of Chapter~\ref{ch:processor}: they are
not among the thirty-two registers that a \riscv programmer sees.

\begin{figure}[htb]
\centering
\begin{tikzpicture}
  \node[blkacc,minimum width=44mm,minimum height=10mm] (ma) at (0,1.5) {\textbf{MA}: memory address register};
  \node[blkamb,minimum width=44mm,minimum height=10mm] (mb) at (0,0) {\textbf{MB}: memory buffer register};
  \node[blk,fill=prgsteel,text=white,minimum width=44mm,minimum height=9mm] (cu) at (0,-1.5) {control unit};
  \node[draw=prgsteel,dashed,rounded corners,fit=(ma)(mb)(cu),inner sep=3mm,label={[lbl]above:Processor}] {};
  \node[blkvio,minimum width=26mm,minimum height=40mm] (m) at (8,0) {\textbf{Memory}};
  \draw[->,bus,prgaccent] (ma.east) -- (ma.east-|m.west) node[midway,above,lbl,text=prgaccent]{address bus};
  \draw[<->,bus,prgamber] (mb.east) -- (mb.east-|m.west) node[midway,above,lbl,text=prgamber]{data bus};
  \draw[->,thick,prgred] (cu.east) -- (cu.east-|m.west) node[midway,above,lbl,text=prgred]{control: read enable, write enable};
\end{tikzpicture}
\caption{The interface between processor and memory.}
\label{fig:mem:iface}
\end{figure}

\subsection{Fetching a word: \texttt{lw t1, 0(t0)}}

\begin{table}[htb]
\centering\small
\caption{The steps of a load.}\label{tab:mem:lw}
\begin{tabularx}{\textwidth}{@{}cLl@{}}
\toprule
\textbf{Step} & \textbf{What happens} & \textbf{Where}\\
\midrule
1 & the address, $\reg{t0}+0$, is computed by the ALU and placed in MA & processor\\
2 & MA drives the address bus & address bus\\
3 & the control unit raises the read enable & control lines\\
4 & the memory decodes the address and puts the four bytes of the word on the data bus & memory, data bus\\
5 & the datum on the data bus is latched in MB & processor\\
6 & MB is copied into the destination register \reg{t1} & processor\\
\bottomrule
\end{tabularx}
\end{table}

A store, \code{sw t1, 0(t0)}, is symmetrical (Table~\ref{tab:mem:lw}): the
address goes into MA and the datum into MB, the write enable is raised, and
the memory stores what it finds on the data bus. The access time is the time
that passes between step 2 and step 5.

\subsection{Synchronising processor and memory}

How does the processor know that the memory has finished? With the
\emph{asynchronous} technique it waits until the memory signals, on an
\term{acknowledge} line, that the operation is complete: the processor adapts
to the speed of the memory, and the design of the memory becomes independent
of that of the processor, at the price of more complexity. With the
\emph{synchronous} technique one of the control lines carries the pulses of a
clock of fixed frequency, and a read or a write completes within one clock
cycle, which must therefore be sized on the speed of the memory. The second
technique is simpler, and is the rule inside a single chip such as a
microcontroller.

% ----------------------------------------------------------------------------
\section{Memory and input/output}

\subsection{Connecting processor, memory and I/O}

The input/output subsystem lets the processor interact with input devices
(keyboard, mouse), output devices (screen, printer) and devices that are both
(mass storage, network cards, serial ports). It can be connected in two ways.

\begin{definition}{Memory-mapped I/O}
With \term{memory-mapped I/O} the memory and the devices share the same bus
and the same address space: a register of a device is read and written like a
memory location, and no special instructions are needed.
\end{definition}

The alternative is \term{isolated I/O}: two physical buses, or two address
spaces, one for memory and one for I/O, and dedicated instructions to reach
the devices. The Intel x86 family has both: the instructions \ins{in} and
\ins{out} besides memory-mapped devices.

\begin{prgmachine}
Memory-mapped I/O is the \riscv way: \riscv has no input/output instructions
at all. Writing a word to a peripheral register with \ins{sw} is how a
\riscv program lights a pixel or reads a button.
\end{prgmachine}

\subsection{Where things live on the ESP32-C6}

Table~\ref{tab:mem:map} is a simplified map of the address space of the chip
of the kit. One address space serves for everything: the same thirty-two
address lines reach ROM, RAM, flash and every device, and the same \ins{lw}
and \ins{sw} operate them. (The table is simplified from the technical
reference manual \citep{esp32c6trm}, which should be consulted before relying
on an address.)

\begin{table}[htb]
\centering\small
\caption{The address space of the ESP32-C6, simplified.}\label{tab:mem:map}
\begin{tabularx}{\textwidth}{@{}lLl@{}}
\toprule
\textbf{Range begins with} & \textbf{What is there} & \textbf{Reached by}\\
\midrule
\hex{4000\ldots} & ROM: the code that starts the chip & instruction fetch\\
\hex{4080\ldots} & SRAM: 512 KB for data, stack and cartridges loaded in RAM & \ins{lw}, \ins{sw}, instruction fetch\\
\hex{4200\ldots} & external flash, seen through a cache: firmware and constants & \ins{lw}, instruction fetch\\
\hex{6000\ldots} & peripheral registers: GPIO, SPI, I2S, timers, Wi-Fi & \ins{lw}, \ins{sw}\\
\bottomrule
\end{tabularx}
\end{table}

The address \hex{40800000}, at which every cartridge is loaded
(Chapter~\ref{ch:setup}), is the beginning of the SRAM.

\subsection{The three component buses, and DMA}

The bus that connects processor, memory and I/O is thus made of three buses:
the \emph{data bus}, which carries the data; the \emph{address bus}, which
carries the address, of memory or of I/O, from or to which the datum is
directed; and the \emph{control lines}: read and write enable, acknowledge,
and others.

\begin{definition}{Direct memory access}
With \term{direct memory access} (DMA), blocks of data go directly between a
device and the memory. The processor starts the transfer and is told when it
is over.
\end{definition}

If several units share the bus and can access the memory directly, a DMA
controller arbitrates among them. On the ESP32-C6 the SPI and I2S peripherals
use DMA to send the frame buffer to the display and the audio samples to the
amplifier: the processor does not copy the pixels one by one.

% ----------------------------------------------------------------------------
% Generated by tools/textbook/export_listings.js from decks/plan.js and code/l15.S: do not edit.
\section{The program of this chapter}\label{sec:lab:l15}

The listings of this chapter are parts of one program, \fname{code/l15.S} in the
repository of the course, printed in full in Listing~\ref{lst:lab:l15}. The
automatic tests of the course material build it and run it on the headless
runner of \prgqt{}, with scripted input, and check the frame;
Figure~\ref{fig:lab:l15} is the frame that the test wrote.

\begin{description}[style=nextline,leftmargin=1.5em]
\item[On the screen.] Five rows of lamps. The first is the word 0x12345678 read with lw. The other four are the bytes at addresses A, A+1, A+2, A+3 read with lbu: 0x78, 0x56, 0x34, 0x12. The lowest address holds the least significant byte: little-endian.
\item[What to press.] Nothing: the program runs by itself.
\end{description}

\begin{figure}[htb]
\centering
\includegraphics[width=0.5\textwidth]{lab_l15.png}
\caption{\fname{l15.S} on \prgqt{}.}\label{fig:lab:l15}
\end{figure}

\asmfile[caption={The complete program of the chapter, \fname{code/l15.S}.},label=lst:lab:l15]{code/l15/program.S}

To run it on \prgqt{} the commands are those of Listing~\ref{lst:lab:l15:run};
Section~\ref{sec:prg:testing} gives the same steps for QEMU and for the board.

\begin{shell}[caption={Building \fname{l15.S} and running it on \prgqt{}.},label=lst:lab:l15:run]
# 1. start PRG32-QT on this computer (it listens on port 8080)

# 2. in the PRG32 folder, with ESP-IDF active: build the cartridge
python3 -m prg32 cartridge build ~/arc12/code/l15.S --portable \
    --entry-prefix l15 --name l15 --out l15.prg32

# 3. send it to PRG32-QT and start it
python3 -m prg32 esp32c6 upload-and-run l15.prg32 \
    --url http://localhost:8080 --slot cart0

# 4. look, press, compare with the previous slide
\end{shell}


% ----------------------------------------------------------------------------
\section{Check your understanding}

\begin{quiz}{A memory has 16 address lines and is byte-addressable. How large is it?}
\item 16 bytes
\item 32 KiB
\item 64 KiB
\item 16 MiB
\end{quiz}

\begin{quiz}{On \riscv, \code{value: .word 0xCAFE0001} is at address $A$. What does \code{lbu t0, 0(t0)} read if $\reg{t0}=A$?}
\item \hex{CA}
\item \hex{FE}
\item \hex{00}
\item \hex{01}
\end{quiz}

\begin{quiz}{Why must DRAM be refreshed?}
\item Because its capacitors lose their charge
\item Because reading is slow
\item Because it is not random access
\item Because the clock is too fast
\end{quiz}

\begin{quiz}{How does a \riscv program write to an output device?}
\item With the \ins{out} instruction
\item With \ins{sw} to the address of a device register
\item With a special register \reg{x32}
\item It cannot: only the firmware can
\end{quiz}

\section{Exercises}

\begin{exercise}[paper first, then \prg]
\begin{enumerate}[nosep]
\item A memory of 4~KiB words of 16 bits: how many address lines, how many
  data lines, how many bits in all?
\item Write the bytes stored, in order of address, by \code{.word 0xDEADBEEF}
  and by \code{.asciz "Ciao"} on a little-endian machine.
\item List the steps of \code{sw t1, 8(t0)} on the processor--memory interface.
\item In \prgqt, find your counter variable in the memory view and watch its
  four bytes change frame by frame.
\item Store a word with \ins{sw} and read it back with four \ins{lbu}
  instructions; rebuild it with shifts and \ins{or}s and check that it
  matches.
\end{enumerate}
\expected{On paper: 12 address lines, 16 data lines, 65\,536 bits;
\texttt{EF BE AD DE} and \texttt{43 69 61 6F 00}. On \prg: the lowest byte of
the counter changes at every frame, the next one every 256 frames. In the
last task the bytes are shifted by 8, 16 and 24 places and then joined.}
\end{exercise}

\section{Glossary}

\begin{glossarytable}
main memory & memoria centrale\\
volatile / persistent & volatile / persistente\\
magnetic core & nucleo magnetico (di ferrite)\\
refresh & rinfresco (refresh)\\
byte order & ordinamento dei byte\\
access time & tempo di accesso\\
aligned word & parola allineata\\
address / data bus & bus indirizzi / bus dati\\
read / write enable & abilitazione lettura / scrittura\\
memory-mapped I/O & I/O mappato in memoria\\
\end{glossarytable}

\section{Conclusions and reading}

Memory is a fundamental component of a computer: a device able to keep the
state of data. Different kinds of memory serve different purposes, because no
technology is fast, large and cheap together. Data are stored as words of
bits, each with an address; \riscv addresses bytes and orders them
little-endian. Access happens through selection (the address bus), reading or
writing (the data bus) and control lines. In \riscv, input and output devices
are reached exactly like memory.

\paragraph{Reading.} \citet{patterson2020}, sections 2.3 and 5.1--5.2;
\citet{stallings2021}, the chapter on internal memory; the technical reference
manual of the ESP32-C6, ``System and Memory'' \citep{esp32c6trm}.
